% Page de résumé en français
\newpage
\thispagestyle{empty}

\begin{center}
    \vspace{2cm}
    {\Large \textbf{\color{blueuniv}RÉSUMÉ}}
    \vspace{1cm}
\end{center}

\noindent
Face à l'augmentation constante des cyberattaques et à l'interconnexion croissante des systèmes informatiques, le développement de systèmes de détection d'intrusions efficaces devient crucial pour la sécurité des infrastructures numériques. Cette recherche propose une approche innovante basée sur l'apprentissage automatique pour améliorer la détection et la classification des intrusions réseau.

\vspace{0.5cm}

\noindent
Notre méthodologie adopte une stratégie hiérarchique en deux étapes appliquée au jeu de données NSL-KDD. La première phase réalise une classification binaire distinguant le trafic normal des activités malveillantes en combinant des algorithmes d'apprentissage supervisé avec un autoencodeur pour la détection d'anomalies. La seconde phase classifie les intrusions détectées en quatre catégories : DoS, Probe, R2L et U2R. Pour traiter le déséquilibre des données, nous intégrons la technique SVM-SMOTE qui génère des échantillons synthétiques pour les classes minoritaires.

\vspace{0.5cm}

\noindent
Les résultats expérimentaux démontrent la supériorité de l'autoencodeur pour la classification binaire avec une précision de 87,52\%. L'application de SVM-SMOTE améliore significativement la détection des attaques sous-représentées, validant l'efficacité de notre approche hiérarchique. Cette recherche contribue au domaine de la cybersécurité en proposant une méthodologie robuste combinant apprentissage supervisé et non-supervisé pour répondre aux défis contemporains de la détection d'intrusions réseau.

\vspace{1cm}

\noindent
\textbf{Mots-clés :} détection d'intrusion, apprentissage automatique, classification hiérarchique, NSL-KDD, autoencodeur, SVM-SMOTE, cybersécurité

\newpage



% Page d'abstract en anglais
\newpage
\thispagestyle{empty}

\begin{center}
    \vspace{2cm}
    {\Large \textbf{\color{blueuniv}ABSTRACT}}
    \vspace{1cm}
\end{center}

\noindent
With the exponential growth of cyber threats and the increasing interconnectivity of computer systems, developing effective intrusion detection mechanisms has become essential for protecting digital infrastructures. This research presents an innovative machine learning-based approach to enhance network intrusion detection and classification capabilities.

\vspace{0.5cm}

\noindent
Our methodology employs a two-stage hierarchical strategy applied to the NSL-KDD dataset. The first stage performs binary classification to distinguish legitimate network traffic from malicious activities by combining supervised learning algorithms with an autoencoder designed for anomaly detection. The second stage categorizes detected intrusions into four attack types: DoS, Probe, R2L, and U2R. To address data imbalance issues, we integrate the SVM-SMOTE technique that generates synthetic samples for minority classes, ensuring better representation of underrepresented attack categories.

\vspace{0.5cm}

\noindent
Experimental results demonstrate the superiority of the autoencoder approach for binary classification, achieving a remarkable accuracy of 87.52\%. The implementation of SVM-SMOTE significantly improves detection rates for underrepresented attacks, validating the effectiveness of our hierarchical framework. This study contributes to the cybersecurity field by proposing a robust methodology that combines supervised and unsupervised learning techniques to address contemporary challenges in network intrusion detection.

\vspace{1cm}

\noindent
\textbf{Keywords:} intrusion detection, machine learning, hierarchical classification, NSL-KDD, autoencoder, SVM-SMOTE, cybersecurity

\newpage